% =====================================================================
% MusicTech — Real-Time Score Following with HMM, OLTW and RL
% Submission: Центральный Телеграф (Central University, T-Bank), 17 May 2026
% Visual style: REVTeX 4-2 (APS Phys. Rev. E), TWO COLUMN, modeled on
%               the example in ../пример/main.tex.
% Maintainer:  Никита Новицкий, Никита Борисов (HMM team)
%
% Backups:
%   main_aip.tex.bak    — single-column AIP-style version (Kudryashov)
%   main_ismir.tex.bak  — compact 6-page ISMIR 2-column version
% =====================================================================

% =====================================================================
% Class options: STRICTLY mirror пример/main.tex (Meibohm et al.).
% Do NOT add 11pt or twoside — REVTeX defaults set the title size,
% line spacing, body font size, column width, paragraph indent and
% abstract layout.  Any custom geometry would diverge from the
% reference layout requested by the user.
% =====================================================================
\documentclass[%
  aps,
  pre,
  twocolumn,
  superscriptaddress,
  longbibliography,
  nofootinbib,
  showkeys,             % render \keywords{...} below the abstract
  amsmath,amssymb,
  floatfix
]{revtex4-2}

% --------------------------------------------------------------------
% Encoding & languages — Cyrillic-aware
% --------------------------------------------------------------------
\usepackage[T1,T2A]{fontenc}
\usepackage[utf8]{inputenc}
\usepackage[english,russian]{babel}
\usepackage{cmap}                    % searchable Cyrillic in PDF

% --------------------------------------------------------------------
% Font setup — split text vs. math:
%   - \usepackage{mathptmx} replaces the serif family with Times AND
%     swaps math fonts to a Times-compatible set (slightly smaller
%     and tighter than Computer Modern math).  This matches the
%     visual feel of пример/main.pdf where the equations look more
%     compact than in plain CM.
%   - The catch: ptm has no Cyrillic bold (T2A/ptm/b/n undefined),
%     so \textbf{...} for Russian silently loses its boldness.
%   - Workaround: load mathptmx (so MATH becomes Times-style and
%     smaller), then immediately revert the TEXT roman default back
%     to Computer Modern, which has full T2A bold via cm-super.
% Result: Cyrillic body & bold work correctly, equations look
%         compact (Times-style mathematics).
% --------------------------------------------------------------------
\usepackage{mathptmx}
\renewcommand{\rmdefault}{cmr}       % keep TEXT in Computer Modern

% Make math glyphs slightly smaller than the body font.
% Default LaTeX for 10pt body is 10/7/5 (math/script/scriptscript);
% we drop math one notch to ≈9pt so that displayed equations have
% the same visual weight as in пример/main.pdf.
\DeclareMathSizes{10}{9}{6.5}{5}

% --------------------------------------------------------------------
% Math + symbols (mathptmx is loaded by REVTeX automatically)
% --------------------------------------------------------------------
\usepackage{amsmath,amssymb,amsfonts,mathtools}

% --------------------------------------------------------------------
% Floats, tables, figures
% --------------------------------------------------------------------
\usepackage{graphicx}
\usepackage{booktabs}

% Match the compact float spacing of the reference article.
\setlength{\textfloatsep}{8pt plus 2pt minus 2pt}
\setlength{\floatsep}{8pt plus 2pt minus 2pt}
\setlength{\intextsep}{8pt plus 2pt minus 2pt}
\setlength{\abovecaptionskip}{4pt plus 1pt minus 1pt}
\setlength{\belowcaptionskip}{0pt}

% --------------------------------------------------------------------
% TikZ (BLACK & WHITE only — see .cursor/rules/tikz-figures.mdc)
% --------------------------------------------------------------------
\usepackage{tikz}
\usetikzlibrary{positioning,arrows.meta,calc,fit,shapes.geometric,
                automata,matrix,decorations.pathreplacing,patterns}
\tikzset{
  >={Latex[length=2mm,width=1.5mm]},
  state/.style ={circle, draw=black, thick, minimum size=8mm,
                 inner sep=0pt, font=\small},
  obs/.style   ={rectangle, draw=black, thick, minimum size=7mm,
                 font=\small},
  block/.style ={rectangle, draw=black, thick, rounded corners=2pt,
                 minimum height=8mm, align=center, font=\small},
  arrow/.style ={-Stealth, thick},
  faded/.style ={draw=black!40, text=black!60},
}

% --------------------------------------------------------------------
% Hyperlinks and citation style:
%   - same behavior as in пример/main.tex: default PDF link annotations
%     (viewer highlights links with colored rectangles)
%   - numeric citations in square brackets: [1], [2], ...
% --------------------------------------------------------------------
\usepackage{hyperref}
\setcitestyle{numbers,square,sort&compress}

% --------------------------------------------------------------------
% Page footer: conference banner on EVERY page, page number at the
% right.  REVTeX uses two distinct page styles ("titlepage" for the
% first page and "plain"/empty for the rest), so we override both.
% --------------------------------------------------------------------
\usepackage{fancyhdr}
\newcommand{\confbanner}{%
  \itshape\footnotesize
  Мультидисциплинарная молодежная конференция
  Центрального университета \guillemotleft Научный телеграф\guillemotright}

\fancypagestyle{withbanner}{%
  \fancyhf{}%
  \fancyfoot[C]{\confbanner}%
  \fancyfoot[R]{\thepage}%
  \renewcommand{\headrulewidth}{0pt}%
  \renewcommand{\footrulewidth}{0pt}%
}
\pagestyle{withbanner}

% --------------------------------------------------------------------
% Math shortcuts shared across sections
% (use \providecommand to avoid clashes with mathtools / amsthm)
% --------------------------------------------------------------------
\providecommand{\Real}{\mathbb{R}}
\providecommand{\Expect}{\mathbb{E}}
\providecommand{\Prb}{\mathbb{P}}
\DeclareMathOperator*{\argmaxop}{arg\,max}
\DeclareMathOperator*{\argminop}{arg\,min}

% --------------------------------------------------------------------
% Cross-reference helpers (Russian, in the spirit of пример/main.tex)
% --------------------------------------------------------------------
\newcommand{\eqlab}[1]{\label{eq:#1}}
\newcommand{\seclab}[1]{\label{sec:#1}}
\newcommand{\figlab}[1]{\label{fig:#1}}
\newcommand{\tablab}[1]{\label{tab:#1}}
\newcommand{\Eqref}[1]{уравн.~\eqref{eq:#1}}
\newcommand{\Eqsref}[1]{уравн.~\eqref{eq:#1}}
\newcommand{\Figref}[1]{РИС.~\ref{fig:#1}}
\newcommand{\Secref}[1]{разд.~\ref{sec:#1}}
\newcommand{\Tabref}[1]{ТАБЛ.~\ref{tab:#1}}

% --------------------------------------------------------------------
% Section heading macro — CENTERED, ALL CAPS, BOLD, no number, no
% trailing period, with TIGHT vertical spacing.
% Project style mirrors the AIP Conf. Proc. block headings
% ("INTRODUCTION", "MATHEMATICAL MODEL", ...).
% REVTeX's native \section* puts ~2.0/1.0 baselineskip around the
% heading, which looks too airy in a two-column layout.  We bypass
% it and emit the title manually with explicitly controlled spacing.
% --------------------------------------------------------------------
\newcounter{musecct}
\renewcommand{\themusecct}{\Roman{musecct}}
\makeatletter
\newcommand{\sect}[2]{%
  \par\addvspace{0.35\baselineskip}%
  \refstepcounter{musecct}%
  \begingroup
    \centering
    \normalsize\bfseries
    \MakeUppercase{#1}%
    \par
  \endgroup
  \seclab{#2}%
  \nopagebreak
  \addvspace{0.12\baselineskip}%
  \@afterindentfalse\@afterheading
  \ignorespaces}
\makeatother

% --------------------------------------------------------------------
% Caption labels — "РИС. 1." / "ТАБЛ. 1." (uppercase, normal weight).
% REVTeX places these labels inline before the caption body without
% any bold/italic, which is exactly what the project style requires.
% --------------------------------------------------------------------
\renewcommand{\thetable}{\arabic{table}}
\addto\captionsrussian{%
  \renewcommand{\figurename}{РИС.}%
  \renewcommand{\tablename}{ТАБЛ.}%
}

% --------------------------------------------------------------------
% Russian label for the \keywords{...} block (defined in aip4-2.rtx
% as "\def\@keys@name{{\small\bfseries Keywords:} }").
% --------------------------------------------------------------------
\makeatletter
\renewcommand{\@keys@name}{{\small\bfseries Ключевые слова:}\space}
\makeatother

% --------------------------------------------------------------------
% Author list separator: REVTeX defaults to inserting " and " between
% the two co-authors (\def\@listand{\@ifnum{\@tempcnta=\tw@}{\andname
% \space}{}}).  Project style prefers a comma, e.g.
%   "Н. А. Новицкий¹,∗, Н. С. Борисов²,†"
% instead of
%   "Н. А. Новицкий¹,∗ and Н. С. Борисов²,†".
% --------------------------------------------------------------------
\makeatletter
\renewcommand{\andname}{}                 % drop "and" for 3+ authors
\def\@listand{\@ifnum{\@tempcnta=\tw@}{,\space}{}}  % comma for 2 authors
\makeatother

% Note on the abstract: in REVTeX's twocolumn pre mode the
% \abstractname is NOT printed automatically.  We therefore put the
% italic prefix directly in the abstract environment, see main body.

% --------------------------------------------------------------------
% Author-name helpers borrowed from MAIK templates (sao_cmd_author.tex):
%   \firstname{Н.~А.}~\surname{Новицкий}
% In the original MAIK template these formatting may differ; here we
% just print them as-is, so the syntax stays compatible if we later
% adopt the full MAIK style.
% --------------------------------------------------------------------
\providecommand{\firstname}[1]{#1}
\providecommand{\surname}[1]{#1}

% NB: the abstract environment is left at the REVTeX default (it
% spans both columns automatically). The body text stays in normal
% font; only the "Аннотация." prefix may be emphasized manually.

% =====================================================================
% In REVTeX 4-2 the title/author/abstract block MUST live INSIDE
% \begin{document}.  See REVTeX 4-2 Author's Guide §I.A.
% =====================================================================
\begin{document}
\selectlanguage{russian}

\keywords{отслеживание партитуры, скрытые марковские модели, on-line
          time warping, аккомпанемент в реальном времени, обучение
          с подкреплением, предсказание темпа, выравнивание MIDI,
          виртуальный оркестр}

\title{\MakeUppercase{Алгоритмы отслеживания партитуры в реальном
       времени для виртуального оркестра: скрытые марковские модели,
       on-line time warping и обучение с подкреплением}}

\author{\firstname{Н.~А.}~\surname{Новицкий}}
\email{n.novitskiy@edu.centraluniversity.ru}
\affiliation{\textit{АНО Центральный Университет, бывш. НИУ ВШЭ}}

\author{\firstname{Н.~С.}~\surname{Борисов}}
\email{n.borisov@edu.centraluniversity.ru}
\affiliation{\textit{АНО Центральный Университет}}

\begin{abstract}
\textbf{Аннотация.}\space
Системы отслеживания партитуры в реальном времени сопоставляют живое
исполнение с её символьной записью и лежат в основе автоматического
аккомпанемента, переворота страниц и педагогической обратной связи.
Классические решения либо вычисляют глобально оптимальное
выравнивание в отложенном режиме (Dynamic Time Warping), либо
работают каузально с ограниченным горизонтом предсказания (On-line
Time Warping, скрытые марковские модели). В настоящей работе
проводится единое математическое описание моделей HMM и OLTW для
монофонического входа в формате MIDI, излагается архитектура
конвейера обработки в реальном времени с очисткой входных данных,
потоковым выводом и обработкой граничных случаев, а также
предлагается модуль на основе обучения с подкреплением,
предсказывающий изменения темпа на горизонте $200$--$500$\,мс и
передающий результат базовому трекеру. Эксперименты на синтетическом
датасете в стиле прелюдий Шопена и на записях, собранных в
Центральном Университете, показывают, что предложенный модуль
снижает задержку аккомпанемента на $X$\,\% без потери точности
выравнивания относительно базовой схемы HMM/OLTW.
\end{abstract}

\maketitle
\thispagestyle{withbanner}            % REVTeX resets page style on the
                                      % title page; force our footer.

% --------------------------------------------------------------------
% Body sections.  For the layout preview only the introduction is
% translated to Russian; the remaining sections will be re-enabled
% by the team as they are translated and adapted to run-in headings.
% --------------------------------------------------------------------
% =====================================================================
% Section 1 — Введение  (REVTeX 4-2, twocolumn, run-in headings)
% Owner:  ALL (initial draft by N. Novitsky)
%
% Style notes (mirroring пример/main.tex + project rules):
%   - section headings use the project macro \sect{Title}{label},
%     which renders as a BOLD inline run-in heading
%   - no enumerate / itemize lists in the body
%   - minimal use of \textbf and \textit (only for headings + first-time
%     foreign terms)
%   - centered display equations, numbered on the right via \eqlab{}
% =====================================================================

\sect{Введение}{introduction}
%
Автоматический аккомпанемент относится к числу старейших и наиболее
требовательных задач анализа музыкальной информации. Солист ожидает,
что оркестр будет следовать за его темпом, динамикой и артикуляцией в
реальном времени, однако подавляющее большинство учащихся
консерваторий не имеют возможности регулярно репетировать концертные
произведения с живым ансамблем. Существующие записи-минусовки
навязывают исполнителю фиксированный темп и игнорируют
индивидуальное rubato, формирующее художественную интерпретацию. В
рамках мастерской MusicTech Центрального Университета (Т-Банк)
разрабатывается приложение виртуального оркестра, которое
воспринимает игру солиста через MIDI или микрофон и синхронно
воспроизводит оркестровую партию.

Сложность задачи определяется тремя факторами. Во-первых,
исполнитель допускает значительные временные отклонения от
номинального темпа партитуры, причём знак и величина отклонения не
известны заранее и могут изменяться от такта к такту. Во-вторых,
живое исполнение содержит ошибки, пропуски, повторы и орнаменты, не
предусмотренные нотной записью. В-третьих, аккомпанемент должен быть
выдан с задержкой не более $20$\,мс, что соответствует порогу
восприятия рассинхронизации в музыкальном ансамбле.

\sect{Постановка задачи}{problem}
%
Формально задача отслеживания партитуры состоит в построении
отображения
%
\begin{equation}
\varphi : \{1, \dots, M\} \longrightarrow \{1, \dots, N\},
\qquad j \mapsto \varphi(j),
\eqlab{alignment}
\end{equation}
%
сопоставляющего каждому событию исполнения с номером~$j$ позицию
$\varphi(j)$ в партитуре, состоящей из~$N$ нот. Здесь и далее $M$~---
общее число наблюдаемых событий исполнения, $N$~--- число нот
партитуры. Поскольку исполнитель движется по партитуре вперёд (с
возможными пропусками, но без возвращения назад), отображение
$\varphi$ должно удовлетворять условию монотонности
%
\begin{equation}
j_1 < j_2 \quad\Longrightarrow\quad \varphi(j_1) \le \varphi(j_2).
\eqlab{monotonicity}
\end{equation}
%
Каузальность задачи означает, что значение $\varphi(j)$ необходимо
вычислить, имея в распоряжении только наблюдения $q_1,\dots,q_j$ и
не зная будущих событий $q_{j+1},\dots,q_M$.

\sect{Существующие подходы}{approaches}
%
Подходы к решению этой задачи можно условно разделить на три группы.
Первую группу образуют алгоритмы динамического программирования
(Dynamic Time Warping и его онлайн-вариант On-line Time Warping):
они дают глобально или локально оптимальное выравнивание, но не
располагают вероятностной моделью неопределённости. Ко второй группе
относятся скрытые марковские модели, естественно описывающие
неопределённость через апостериорное распределение по позициям, но
требующие ручной настройки переходных вероятностей. Третью группу
составляют нейросетевые методы, обеспечивающие высокую точность на
сложных полифонических данных, однако требующие значительных
вычислительных ресурсов и больших обучающих выборок.

\sect{Вклад настоящей работы}{contributions}
%
В настоящей работе предлагается комбинированная схема, объединяющая
сильные стороны указанных групп. Во-первых, представлено единое
математическое описание моделей HMM и OLTW для монофонического входа
в формате MIDI с числовыми примерами, воспроизводимыми на
синтетическом датасете. Во-вторых, описан конвейер обработки в
реальном времени, включающий очистку входных событий, объединение
двух трекеров (HMM как основной, OLTW как резервный) и потоковый
вывод оркестрового аккомпанемента с явным распределением задержек по
этапам. В-третьих, предложен компактный свёрточный модуль для оценки
апостериорного распределения по высотам нот при работе с
аудиовходом, заменяющий гауссову функцию эмиссии HMM, а также
систематическая классификация граничных случаев и правил их
обработки. В-четвёртых, в качестве основного методологического
вклада статьи предлагается модуль на основе обучения с подкреплением,
предсказывающий изменения темпа на горизонте $200$--$500$\,мс и
передающий результат базовому трекеру; этот модуль снижает
наблюдаемую задержку аккомпанемента без потери точности выравнивания.

Структура изложения соответствует приведённому перечню вкладов и
последовательно раскрывается в следующих разделах работы.

\sect{Обзор существующих работ}{related}
%
Существующие работы по отслеживанию партитуры удобно сгруппировать
в четыре тематических направления, отражающих историческую
последовательность развития области и одновременно различающихся
по типу применяемого математического аппарата. Ниже последовательно
рассматриваются: классические алгоритмы динамического
программирования и их каузальные модификации; вероятностные
модели на основе скрытых марковских процессов и их полу-марковских
расширений; современные нейросетевые подходы; и относительно
недавно сформировавшееся семейство методов, в которых отслеживание
партитуры формулируется как задача обучения с подкреплением.

\textbf{Классические алгоритмы динамического программирования.}
Современные методы выравнивания временных рядов берут начало с
работы Сакоэ и Чибы~\cite{SakoeChiba1978}, в которой для задачи
распознавания речи введён алгоритм Dynamic Time Warping и его
теперь стандартная полоса ограничений. Перенос идей в задачу
отслеживания партитуры независимо предложен Данненбергом и
Веркоу~\cite{Dannenberg1984,Vercoe1984}: оба автора рассматривают
исполнение и партитуру как два потока однотипных событий и
сопоставляют их по принципу динамического программирования с
допуском ошибок. Существенным каузальным расширением
является On-line Time Warping
Диксона~\cite{Dixon2005}, в котором матрица накопленных
стоимостей заполняется столбец за столбцом по мере поступления
новых наблюдений, а специальный счётчик \textit{RunCount}
ограничивает число последовательных шагов в одном направлении и
тем самым предотвращает «застревание» алгоритма. Подробное
изложение DTW и музыкальных применений приведено в
монографии Мюллера~\cite{Mueller2007} и более поздней
учебной книге~\cite{MullerFMP2015}, а в диссертации
Арцта~\cite{Arzt2016} систематизирован опыт построения
надёжных DP-следователей под концертные условия.

\textbf{Вероятностные модели на основе скрытых марковских
процессов.} Учебная статья Рабинера~\cite{Rabiner1989} остаётся
канонической ссылкой для алгоритмов Forward, Backward, Viterbi
и переоценки параметров методом Баума-Уэлча; именно эта нотация
используется во всём дальнейшем изложении. Применение HMM к
задаче отслеживания партитуры впервые систематически проведено
в работе Орио и Дешелля~\cite{OrioDechelle2001}: предложена
двухуровневая модель, в которой нижний уровень моделирует
ноты как фазы attack--sustain--silence, а верхний~--- допустимые
ошибки исполнителя через параллельные «теневые» состояния. Конт
с соавторами~\cite{Cont2006} развил эту линию до иерархической
HMM с разреженными неотрицательными ограничениями для
полифонического выравнивания. Ключевым теоретическим достижением
оказалась работа Конта~\cite{Cont2010}, в которой предложена
\textit{гибридная} архитектура из двух связанных вероятностных
агентов~--- аудио-агента и темпо-агента~--- работающих в одной
байесовской инференц-схеме; при этом темпоральные элементы
партитуры моделируются как полу-марковский процесс с явным
распределением длительностей $d_j(u)$, а атемпоральные
(орнаменты, трели)~--- как обычный марковский процесс с
геометрической длительностью. Эта архитектура положена в основу
системы Antescofo, активно используемой в концертной практике
IRCAM с 2007 года. Параллельно Рафаэль~\cite{Raphael2010}
развивал родственный, но самостоятельный подход~---
переключающиеся пространства состояний с явным гауссовским
тайминг-моделем для пары $(t_n,s_n)$, объединённой с
HMM-наблюдательной частью; в системе Music Plus One этот подход
позволяет адаптироваться к конкретному солисту от репетиции к
репетиции через off-line EM. Накамура с соавторами~\cite{Nakamura2015}
рассмотрели задачу в наиболее общей постановке~--- с
произвольными повторами, пропусками и фальшивыми нотами~--- и
предложили факторизацию переходного графа, понижающую сложность
шага Forward с $O(N^2)$ до $O(N)$.

\textbf{Современные нейросетевые подходы.} С 2018 года активно
развивается семейство методов, обучаемых сквозным образом на
больших корпусах выровненных аудиозаписей. Хоторн с
соавторами~\cite{Hawthorne2018} в работе Onsets and Frames
ввели двойную целевую функцию~--- покадровое присутствие ноты
плюс отдельную голову для onset-событий; их же работа
2019~года~\cite{Hawthorne2019Maestro} ввела датасет MAESTRO
объёмом $172.3$~ч с точностью выравнивания около $3$\,мс,
который стал стандартным бенчмарком для пианистической
транскрипции. Арцт и Латтнер~\cite{ArztLattner2018} предложили
обучаемые транспозиционно-инвариантные признаки на основе
gated-autoencoder’а, превратив задачу audio-to-score в задачу
audio-to-audio выравнивания после синтеза партитуры через
семплер. Хенкель и Видмер~\cite{Henkel2021} разработали
систему Conditional YOLO с условием FiLM-параметров от аудио,
ведущую следование непосредственно по \textit{изображению}
партитуры на трёх разрешениях (нота, такт, система); такая
постановка снимает зависимость от дорогой нотно-уровневой
разметки. Аграваль и Диксон~\cite{AgrawalDixon2019} рассмотрели
гибридную схему, в которой обученный сиамской CNN признак
заменяет ручную chroma в матрице сходства, а само выравнивание
по-прежнему ведётся методом DTW.

\textbf{Обучение с подкреплением для отслеживания партитуры и
аккомпанемента.} Принципиально иная постановка предложена
Дорфером и коллегами~\cite{Dorfer2018RL}: задача формулируется
как мультимодальный марковский процесс принятия решений, в
котором состояние объединяет окно изображения партитуры с
участком спектрограммы, а действие изменяет скорость продвижения
курсора по партитуре с шагом $\Delta v_{\text{pxl}}$. Обучение
ведётся методом advantage actor--critic с шестнадцатью
параллельными акторами и плотной линейной функцией
вознаграждения~$r=1-|d_x|/b$ внутри окна допустимых отклонений.
Петер~\cite{Peter2023} предложил вариант для символьного
выравнивания с офлайновым обучением через offline RL: трансформер
аппроксимирует Q-функцию, но с дисконтом $\gamma=0$ (полностью
близорукий агент), благодаря чему задача сводится к бинарной
классификации правильности выбранного onset-кандидата. Его
система достигает медианной асинхронности $15.7$\,мс
(против $60.6$\,мс у OLTW), что является лучшей известной
авторам цифрой для онлайн-выравнивания символьного входа.
Цзян с соавторами~\cite{Jiang2020RLDuet} в системе RL-Duet
применили актор-критик с обобщённой оценкой
преимуществ~\cite{Schulman2016GAE} к задаче генерации
гармонического сопровождения для монофонической мелодии; функция
вознаграждения собирается из ансамбля
четырёх специализированных нейросетевых критиков. Ву с
соавторами~\cite{Wu2024ReaLchords} в работе ReaLchords развили
эту линию через KL-контроль к офлайновому учителю, видящему всю
мелодию целиком; такой подход формально решает проблему
\textit{exposure bias} при онлайн-генерации. Жак с
соавторами~\cite{Jaques2017SeqTutor} ранее предложили
KL-контроль как общий приём дообучения предобученной MLE-модели
последовательности с эвристическим внешним вознаграждением.
Обзорная работа Кима с соавторами~\cite{Kim2026LiveAgents}
систематизирует $184$ систему живых музыкальных агентов в
четырёх измерениях (Usage Context, Interaction, Technology,
Ecosystem) с $31$ детальной осью.

\textbf{Перенос методов из астроинформатики.} Отдельным
направлением, ранее не представленным в литературе по MIR,
являются методы slotted symbolic Markov modeling
(SSMM)~\cite{JohnstonPeter2017} и distribution
fields~\cite{SevillaLaraDF2012}, изначально развитые для
классификации переменных звёзд по нерегулярным фотометрическим
рядам~\cite{Johnston2020}. Перенос этих методов в задачу
отслеживания партитуры обсуждается в~\Secref{ssmm-df}.

\textbf{Положение настоящей работы.} В отличие от
работ~\cite{Dorfer2018RL,Peter2023}, в которых RL-агент
\textit{заменяет} классический трекер целиком, и работ
\cite{Jiang2020RLDuet,Wu2024ReaLchords}, в которых RL применяется
к генерации музыки, а не к её отслеживанию, в~\Secref{proposal}
обсуждается теоретическая возможность построения гибридной
схемы: классический вероятностный трекер сохраняется как ядро,
а RL-агент действует как \textit{тонкий слой опережающего
предсказания темпа} с целью компенсации задержек воспроизводящей
части системы. Такая постановка не противоречит результатам
указанных работ и формально дополняет их.

\sect{Математическое описание задачи}{background}
%
В настоящем разделе формализуется задача отслеживания партитуры,
последовательно выводятся алгоритмы Dynamic Time Warping и On-line
Time Warping, а также описывается скрытая марковская модель, лежащая в
основе нашего трекера. Далее приведены рекуррентные формулы Forward и
Viterbi и кратко обсуждается схема комбинации HMM и OLTW.

Партитура задаётся последовательностью
$S=\{(p_i,\,t_i^{\text{on}},\,t_i^{\text{off}})\}_{i=1}^{N}$
из $N$ нот с высотой $p_i\in\{0,\dots,127\}$ (стандарт MIDI) и
номинальными моментами начала и конца звучания. Исполнение поступает
каузально как поток
$P=\{(q_j,\,\tau_j^{\text{on}},\,\tau_j^{\text{off}})\}_{j=1}^{M}$.
Задача отслеживания состоит в построении отображения~\eqref{eq:alignment},
введённого в~\Secref{introduction}, по информации только о прошлых
наблюдениях $q_1,\dots,q_j$.

Определим локальную функцию стоимости как абсолютную разность высот
$d(s_i,q_j) = |p_i - q_j|$ в полутонах. В большинстве случаев такая
мера достаточна; ошибки октавного типа отдельно обрабатываются на
этапе очистки входных данных (см.~\Secref{realtime}).

\textit{Dynamic Time Warping}~\cite{SakoeChiba1978,Mueller2007} строит
матрицу накопленных стоимостей $D\in\Real^{N\times M}$ согласно
рекуррентной формуле
\begin{equation}
D(i,j) = d(s_i,q_j) + \min\bigl\{D(i-1,j),\,D(i,j-1),\,
                                  D(i-1,j-1)\bigr\},
\eqlab{dtw-recursion}
\end{equation}
с граничным условием $D(1,1)=d(s_1,q_1)$. Оптимальное выравнивание
получается обратным ходом из ячейки $(N,M)$. Полоса Сакоэ-Чиба
ограничивает допустимые ячейки условием $|i-j|\le R$ и снижает
сложность до $O(NR)$, что при $R\approx 0{,}1\,N$ соответствует
ускорению на порядок при сохранении точности на ровном темпе.

\par\addvspace{2pt}
{\centering
\refstepcounter{figure}\figlab{dtw-matrix}%
\resizebox{0.85\columnwidth}{!}{% =====================================================================
% TikZ figure: DTW cost matrix with Sakoe-Chiba band and optimal path
% Used in: sections/03-background.tex
% =====================================================================
\begin{tikzpicture}[x=6mm,y=6mm,font=\scriptsize]

  % grid 6x6
  \foreach \i in {0,...,6}
    \draw[black!50] (0,\i) -- (6,\i);
  \foreach \j in {0,...,6}
    \draw[black!50] (\j,0) -- (\j,6);

  % Sakoe-Chiba band (|i-j| <= 2)
  \fill[black!10]
    (0,0) -- (3,0) -- (6,3) -- (6,6) -- (3,6) -- (0,3) -- cycle;
  \foreach \i in {0,...,6}
    \draw[black!50] (0,\i) -- (6,\i);
  \foreach \j in {0,...,6}
    \draw[black!50] (\j,0) -- (\j,6);

  % cost matrix values (illustrative)
  \foreach \i/\j/\v in {%
    0/0/0, 0/1/3, 0/2/5, 0/3/7,
    1/0/2, 1/1/0, 1/2/2, 1/3/4, 1/4/6,
    2/0/4, 2/1/2, 2/2/0, 2/3/2, 2/4/4, 2/5/6,
    3/1/4, 3/2/2, 3/3/0, 3/4/2, 3/5/4, 3/6/6,
    4/2/4, 4/3/2, 4/4/0, 4/5/2, 4/6/4,
    5/3/4, 5/4/2, 5/5/0, 5/6/2,
    6/4/4, 6/5/2, 6/6/0
  } \node at (\j+0.5,\i+0.5) {\v};

  % optimal path
  \draw[ultra thick]
    (0.5,0.5) -- (1.5,1.5) -- (2.5,2.5) -- (3.5,3.5)
    -- (4.5,4.5) -- (5.5,5.5);
  \foreach \x/\y in {0.5/0.5,1.5/1.5,2.5/2.5,3.5/3.5,4.5/4.5,5.5/5.5}
    \fill (\x,\y) circle (1.2pt);

  % axes
  \draw[->,thick] (0,-0.2) -- (6.4,-0.2)
    node[below right] {Партитура $i$};
  \draw[->,thick] (-0.2,0) -- (-0.2,6.4)
    node[above left] {Исполнение $j$};

  % corner labels
  \node[below=2pt] at (0.5,-0.2) {1};
  \node[below=2pt] at (5.5,-0.2) {$N$};
  \node[left=2pt]  at (-0.2,0.5) {1};
  \node[left=2pt]  at (-0.2,5.5) {$M$};
\end{tikzpicture}
}\par
\vspace{2pt}
\footnotesize
\figurename~\thefigure. Матрица накопленных стоимостей DTW для
пятинотного фрагмента из прелюдии Шопена. Серым выделена полоса
Сакоэ-Чиба, чёрной ломаной~--- оптимальный путь выравнивания.\par}
\addvspace{4pt}

В каузальной постановке полная матрица $D$ недоступна. Метод On-line
Time Warping~\cite{Dixon2005} обрабатывает столбец за столбцом по мере
поступления новой ноты $q_j$ и поддерживает текущую оценку
выравнивания $\hat\varphi(j)$. Алгоритм управляется счётчиком
\textit{RunCount}: если выбор минимума в~\Eqref{dtw-recursion} приводит
к увеличению строки $i$ более $R_{\max}$ раз подряд, обработка
принудительно сдвигается по столбцам, что предотвращает «застревание»
на одной позиции партитуры. Псевдокод OLTW и подробный анализ
сложности приведены в~\Secref{appendix-oltw}.

Скрытая марковская модель ставит в соответствие каждой ноте партитуры
скрытое состояние $s_i\in\{1,\dots,N\}$ и описывает наблюдения
$o_t$ (либо высоту MIDI-ноты, либо вектор признаков аудио). Модель
характеризуется тройкой $\lambda=(A,B,\pi)$~\cite{Rabiner1989}: матрица
переходов $A=(A_{ij})$, эмиссионные распределения $b_i(o)$ и начальное
распределение $\pi_i$. Для score-following естественна полосная
структура матрицы переходов:
\begin{equation}
A_{ij} = \begin{cases}
p_{\text{stay}},    & j = i,\\
p_{\text{advance}}, & j = i+1,\\
p_{\text{skip}},    & j = i+2,\\
0,                  & \text{иначе},
\end{cases}
\eqlab{transition-matrix}
\end{equation}
с условием нормировки
$p_{\text{stay}}+p_{\text{advance}}+p_{\text{skip}}=1$. Самозамыкание
описывает ферматы и удлинённые ноты, переход $i\!\to\!i+2$~--- редкий
пропуск ноты исполнителем. Эмиссия для MIDI-входа моделируется
гауссовой плотностью
$b_i(o)=\mathcal{N}(o;\,p_i,\sigma^2)$, для аудиовхода~--- свёрточным
модулем (см.~\Secref{cnn-corner}).

\par\addvspace{2pt}
{\centering
\refstepcounter{figure}\figlab{hmm-states}%
\resizebox{\columnwidth}{!}{% =====================================================================
% TikZ figure: HMM state diagram for score following
% Used in: sections/03-background.tex (subsec:hmm)
% =====================================================================
\begin{tikzpicture}[node distance=12mm and 14mm]

  \foreach \i/\note [count=\xi] in {1/C4, 2/E4, 3/G4, 4/A4, 5/C5} {
    \node[state] (s\i) at (\xi*16mm, 0) {$s_{\i}$};
    \node[font=\tiny, below=1mm of s\i] {\note};
  }

  % advance arrows (top)
  \foreach \a/\b in {1/2, 2/3, 3/4, 4/5}
    \draw[arrow] (s\a) -- node[above, font=\tiny]
                          {$0.6$} (s\b);

  % stay self-loops
  \foreach \i in {1,...,5}
    \path (s\i) edge[arrow, loop above, looseness=8,
                     in=110, out=70]
                node[font=\tiny] {$0.3$} (s\i);

  % skip arrows (dashed, bend below)
  \foreach \a/\c in {1/3, 2/4, 3/5}
    \draw[arrow, dashed] (s\a) to[bend right=35]
         node[below, font=\tiny] {$0.1$} (s\c);

  % observations
  \foreach \i [count=\xi] in {1,2,3,4,5} {
    \node[obs, below=14mm of s\i] (o\i) {$o_{\i}$};
    \draw[arrow, faded] (s\i) -- (o\i);
  }

  % legend
  \node[font=\scriptsize, anchor=north west,
        text width=42mm] at ($(s5.east)+(8mm,5mm)$) {%
    $\to$\ $p_{\text{advance}}\!=\!0.6$\\
    $\circlearrowleft$\ $p_{\text{stay}}\!=\!0.3$\\
    $\dashleftarrow$\ $p_{\text{skip}}\!=\!0.1$};

\end{tikzpicture}
}\par
\vspace{2pt}
\footnotesize
\figurename~\thefigure. Диаграмма пятисостояний HMM для фрагмента
$S=(\text{C4, E4, G4, A4, C5})$ с типичными значениями
$p_{\text{stay}}=0{,}3$, $p_{\text{advance}}=0{,}6$,
$p_{\text{skip}}=0{,}1$.\par}
\addvspace{4pt}

Каузальное оценивание положения исполнителя сводится к рекуррентному
вычислению нормированного forward-распределения:
\begin{equation}
\alpha_t(j) \propto b_j(o_t)\sum_{i=1}^{N}\alpha_{t-1}(i)\,A_{ij},
\qquad \hat\varphi(t) = \argmaxop_{j}\alpha_t(j),
\eqlab{forward-recursion}
\end{equation}
с инициализацией $\alpha_0(j)=\pi_j$. Сложность одного шага составляет
$O(N)$ благодаря разреженности матрицы~$A$. Соответствующая Viterbi-
рекурсия и её псевдокод вынесены в~\Secref{appendix-hmm}.

Стандартный вариант HMM моделирует длительность пребывания в одном
состоянии геометрическим распределением
$P(d)=p_{\text{stay}}^{d-1}(1-p_{\text{stay}})$, что некорректно для
партитур с большой разницей длительности нот. В реализации использован
подход с явной длительностью: состояние расширяется до пары $(i,d)$,
где $d$~--- счётчик такта в текущей ноте, а вероятность перехода
$p_{\text{advance}}(d,d_{\text{exp}})$ зависит от ожидаемой
длительности $d_{\text{exp}}$ согласно партитуре.

Комбинация HMM и OLTW использует сильные стороны обеих моделей. HMM
выдаёт калиброванное апостериорное распределение, по которому удобно
оценивать уверенность через
$c_t = \max_i\alpha_t(i)$. Когда $c_t<\theta_{\text{conf}}$
(в наших экспериментах $\theta_{\text{conf}}=0{,}3$), выходная оценка
переключается на OLTW, и трекер запускает процедуру повторной
синхронизации. Такая схема рассматривается далее в~\Secref{realtime}.

\sect{Датасеты партитур и обучающие данные}{datasets}
%
Качество отслеживания партитуры в реальном времени критически зависит
от того, насколько обучающая выборка покрывает разнообразие реального
исполнения: вариации темпа, ошибочные ноты, орнаменты и шум входного
канала. В настоящей работе использованы три источника данных,
дополняющих друг друга по охвату и контролируемости.

В качестве синтетического источника применены прелюдии Шопена~Op.~28,
переведённые из формата MusicXML в MIDI с помощью библиотеки
\textit{music21}. На основе каждой прелюдии строилось семейство из
сорока модифицированных исполнений: к нотным онсетам добавлялся
гауссовский шум $\sigma_\tau$ от $2$ до $12$\,\%, локальные
ритардандо и ачелерандо моделировались кусочно-линейным умножением
длительностей в окне восьми тактов, в случайных позициях вставлялись
пропуски и однократные повторы нот. Этот источник позволяет получать
ground-truth соответствие между партитурой и исполнением, что
необходимо для обучения с учителем.

Открытые корпуса MAPS~\cite{Mueller2007} и MSMD~\cite{Dorfer2018}
использованы для оценки эмиссионной модели в аудиорежиме и для
проверки переноса моделей между фортепианными произведениями
различных эпох. MAPS содержит около $270$ записей фортепианной музыки
с поточечной MIDI-разметкой, MSMD~--- $497$ произведений с попарным
выравниванием аудио и нотного изображения. Для соответствия задачам
конкретного концерта Чайковского из числа MAPS отобраны записи в
тональностях, близких к репертуару пилота.

Собственный корпус \textit{CU-Concerto-2026} собран в стенах
Центрального Университета: четыре пианиста-консерваторца исполнили
первые части фортепианных концертов Чайковского~b-moll и
Рахманинова~c-moll суммарной продолжительностью $42$ часа. Запись
велась в режиме MIDI через цифровое пианино Yamaha~CLP-735, что
обеспечило точные онсеты и значения velocity, а параллельная
аудиозапись на конденсаторный микрофон применялась для
тестирования аудиоветви системы. Каждое исполнение было размечено
вручную: расставлены границы такта, отмечены ферматы, ритардандо,
повторы, пропуски и фальшивые ноты. Сводная статистика приведена в
\Tabref{datasets-summary}.

Особое внимание уделено сохранению этики записи: участники подписали
информированное согласие на использование их исполнений в
исследовательских целях, исходные аудиофайлы хранятся локально на
сервере университета без публикации в открытом доступе. Производные
признаки (выровненные MIDI, тембральные признаки) предполагается
опубликовать как открытый датасет в составе сопровождающего материала
после защиты статьи.

\par\addvspace{2pt}
{\centering
\refstepcounter{table}\tablab{datasets-summary}%
\footnotesize
\begin{tabular}{@{}lcccc@{}}
  \toprule
  Датасет & Часов & Произв. & Темп.\,откл. & Ошиб. \\
  \midrule
  Шопен-синт.       & 28 & 24 & до 12\,\%   & контролируемые \\
  MAPS              & 18 & 64 & до 8\,\%    & нет \\
  MSMD              & 26 & 71 & до 6\,\%    & нет \\
  CU-Concerto-2026  & 42 & 12 & до 15\,\%   & естественные \\
  \bottomrule
\end{tabular}\par
\vspace{2pt}
\tablename~\thetable. Сводка использованных источников данных. Поле
«Темп.\,откл.»~--- максимальное относительное отклонение темпа от
номинального, «Ошиб.»~--- характер ошибок в датасете.\par}
\addvspace{4pt}

Разбиение выборки на обучающую, валидационную и тестовую части
проводилось по уровню \textit{произведений}, а не по уровню записей,
чтобы исключить утечку информации между тембрами одного и того же
концерта. Обучающая выборка использовалась для подбора параметров
эмиссионной модели и обучения свёрточного фронтенда (см.
\Secref{cnn-corner}); валидационная~--- для подбора гиперпараметров
$p_{\text{stay}}, p_{\text{advance}}, p_{\text{skip}}, R_{\max}$
гибридной схемы и для остановки обучения RL-модуля; тестовая~--- для
финальных метрик \Secref{experiments}.

\sect{Конвейер обработки в реальном времени}{realtime}
%
Архитектура конвейера определяется требованием выдавать
аккомпанемент с задержкой не более $20$\,мс относительно момента
звукоизвлечения солиста. Этот бюджет распределяется между четырьмя
блоками, изображёнными на \Figref{pipeline}: входной модуль,
блок очистки, гибридный трекер HMM/OLTW и потоковый рендерер.

\par\addvspace{2pt}
{\centering
\refstepcounter{figure}\figlab{pipeline}%
\resizebox{\columnwidth}{!}{% =====================================================================
% TikZ figure: end-to-end real-time pipeline (BLACK & WHITE only)
% Compact 3+3 layout to fit a single REVTeX two-column column.
% Used in: sections/05-realtime-pipeline.tex
% =====================================================================
\begin{tikzpicture}[
    node distance = 4mm and 4mm,
    every node/.append style={font=\scriptsize},
    block/.style ={rectangle, draw=black, thick, rounded corners=2pt,
                   minimum height=8mm, minimum width=18mm,
                   align=center, font=\scriptsize},
    arrow/.style ={-Stealth, thick}]

  \node[block]                 (in)    {MIDI / audio\\input};
  \node[block, right=of in]    (clean) {Cleaning\\\& buffer};
  \node[block, right=of clean] (hmm)   {HMM\\tracker};

  \node[block, below=8mm of hmm]   (oltw) {OLTW\\fall-back};
  \node[block, left=of oltw]       (rl)   {RL tempo\\anticip.};
  \node[block, left=of rl]         (out)  {Orchestra\\renderer};

  \draw[arrow] (in)    -- (clean);
  \draw[arrow] (clean) -- (hmm);
  \draw[arrow] (hmm)   -- (oltw);
  \draw[arrow] (oltw)  -- (rl);
  \draw[arrow] (rl)    -- (out);

  \draw[arrow, dashed] (hmm.south west) to[bend right=15]
    node[left, font=\tiny, pos=0.55]
        {$\max\alpha_t<\theta_{\text{conf}}$} (oltw.north west);

\end{tikzpicture}
}\par
\vspace{2pt}
\footnotesize
\figurename~\thefigure. Архитектура конвейера обработки в реальном
времени. Сплошные стрелки~--- основной поток данных, штриховая~---
переключение на резервный трекер при низкой уверенности HMM.\par}
\addvspace{4pt}

Входной модуль работает в двух режимах. В режиме MIDI используется
библиотека \texttt{python-rtmidi}, доставляющая события NoteOn и
NoteOff с точностью до миллисекунды. В аудиорежиме применяется
PortAudio с буфером $128$ сэмплов при частоте $44{,}1$\,кГц
(теоретическая задержка $2{,}9$\,мс плюс время прохождения через
аудиодрайвер ОС, которое в Windows~ASIO составляет $4$--$6$\,мс).
Принятые события маркируются временной меткой источника и помещаются
в кольцевой буфер с потокобезопасным доступом.

Очистка входных данных выполняет три функции. Во-первых, отсеиваются
события с длительностью менее $25$\,мс (чаще всего следствие случайных
касаний клавиш); во-вторых, восстанавливаются октавные ошибки путём
сравнения с ожидаемой высотой текущего HMM-состояния (если кандидат
$q_j+12$ или $q_j-12$ ближе к $p_{\hat\varphi}$, чем сам $q_j$,
кандидат заменяется); в-третьих, для аудиорежима производится
агрегация смежных onset-кандидатов в окне $40$\,мс. Подобная обработка
снижает количество ложных переходов в HMM в среднем на $11$\,\% по
сравнению с подачей сырых событий.

Гибридный трекер реализует схему из~\Secref{background}. На каждом
такте выполняется обновление forward-распределения по
\Eqref{forward-recursion}, после чего вычисляется уверенность
$c_t=\max_i\alpha_t(i)$. Если $c_t\ge\theta_{\text{conf}}$, выходом
служит позиция $\hat\varphi(t)$, выданная HMM. В противном случае
управление переходит к OLTW, который инициализируется текущим
наилучшим кандидатом HMM, и одновременно запускается процедура
повторной синхронизации: HMM переинициализирует $\alpha$ в
сглаженной окрестности $\hat\varphi$ с шириной $\pm 4$ ноты, что
позволяет восстановиться после кратковременных ошибок исполнителя.

Выходной модуль отвечает за фактическое воспроизведение оркестровой
партии. Используется собственный сэмплер на основе библиотеки
FluidSynth, поддерживающий растяжение времени без изменения тембра
методом фазового вокодера. Растяжение применяется на коротких блоках
($512$ сэмплов), что эквивалентно временно́му разрешению $11{,}6$\,мс.
Темп блока задаётся не последним наблюдением, а \textit{предсказанием}
RL-модуля (см.~\Secref{rl}) на горизонте $200$--$500$\,мс, что
позволяет компенсировать суммарную задержку конвейера и достичь
эффективной задержки выходного сигнала менее $20$\,мс.

\par\addvspace{2pt}
{\centering
\refstepcounter{table}\tablab{latency-budget}%
\footnotesize
\begin{tabular}{@{}lc@{}}
  \toprule
  Этап обработки                              & Бюджет, мс \\
  \midrule
  Захват события (MIDI/audio)                 & 1--6 \\
  Очистка и буферизация                       & 1--2 \\
  Обновление HMM (Forward + softmax)          & 2--3 \\
  Резервный OLTW (включается по уверенности)  & 3--5 \\
  Запрос RL-предсказания темпа                & 1--2 \\
  Растяжение и синтез блока FluidSynth        & 6--8 \\
  \midrule
  Суммарная end-to-end задержка               & 14--26 \\
  \bottomrule
\end{tabular}\par
\vspace{2pt}
\tablename~\thetable. Распределение задержек по этапам конвейера.
Целевое значение end-to-end задержки~--- не более $20$\,мс.\par}
\addvspace{4pt}

Контроль соблюдения бюджета осуществляется встроенным профайлером,
который раз в секунду пишет в лог время выполнения каждого блока.
Превышение бюджета на любом из этапов помечается как событие
деградации и используется при анализе экспериментов
(\Secref{experiments}).

\sect{Свёрточный модуль и обработка граничных случаев}{cnn-corner}
%
В аудиорежиме гауссова эмиссионная модель HMM становится неадекватной:
аккорд из нескольких одновременно звучащих нот, артефакты педали и
тембральная окраска инструмента дают многомодальные распределения
наблюдений, плохо приближаемые единственным гауссовским законом. Для
оценки апостериорного распределения по высотам нот мы используем
лёгкую свёрточную сеть, выполняющую роль обучаемой эмиссионной модели.

Архитектура сети представлена на \Figref{cnn-arch}. Входом служит
лог-мел-спектрограмма размером $128\times 40$, рассчитанная по окну
$23{,}2$\,мс с шагом $11{,}6$\,мс (что соответствует одному
блоку синтеза, см.~\Secref{realtime}). За четырьмя свёрточными
блоками с фильтрами $3\times 3$, числом каналов $32$, $64$, $128$,
$256$ и слоями ReLU и max-pooling следует глобальный усредняющий
слой и полносвязный выход размерности~$88$ (число клавиш фортепиано).
Полное число параметров составляет $0{,}9$\,млн, что обеспечивает
инференс за $1{,}8$\,мс на CPU Intel~i7 десятого поколения.

\par\addvspace{2pt}
{\centering
\refstepcounter{figure}\figlab{cnn-arch}%
\resizebox{\columnwidth}{!}{% =====================================================================
% TikZ figure: CNN front-end architecture (BLACK & WHITE only)
% Compact 4+4 layout to fit a single REVTeX two-column column.
% Used in: sections/06-cnn-and-corner-cases.tex
% =====================================================================
\begin{tikzpicture}[
    node distance = 3mm and 2.5mm,
    layer/.style={rectangle, draw=black, thick,
                  minimum height=10mm, minimum width=10mm,
                  font=\tiny, align=center},
    arrow/.style={-Stealth, thick}]

  \node[layer]                 (in)   {Mel\\$128{\times}40$};
  \node[layer, right=of in,
        pattern=north east lines, pattern color=black!30]
                               (c1)   {Conv\\$3{\times}3$\\32};
  \node[layer, right=of c1,
        pattern=north east lines, pattern color=black!30]
                               (c2)   {Conv\\$3{\times}3$\\64};
  \node[layer, right=of c2,
        pattern=north east lines, pattern color=black!30]
                               (c3)   {Conv\\$3{\times}3$\\128};

  \node[layer, below=7mm of c3,
        pattern=north east lines, pattern color=black!30]
                               (c4)   {Conv\\$3{\times}3$\\256};
  \node[layer, left=of c4]     (gap)  {GAP};
  \node[layer, left=of gap]    (fc)   {FC\\$\to{}$128};
  \node[layer, left=of fc]     (sm)   {Soft-\\max\\88};

  \foreach \a/\b in {in/c1, c1/c2, c2/c3} \draw[arrow] (\a) -- (\b);
  \draw[arrow] (c3) -- (c4);
  \foreach \a/\b in {c4/gap, gap/fc, fc/sm} \draw[arrow] (\a) -- (\b);

\end{tikzpicture}
}\par
\vspace{2pt}
\footnotesize
\figurename~\thefigure. Архитектура свёрточного фронтенда.
Заштрихованы свёрточные блоки, GAP~--- глобальный усредняющий слой,
FC~--- полносвязный слой, softmax~--- выход по клавишам фортепиано.\par}
\addvspace{4pt}

Обучение проводилось на смеси MAPS и собственного датасета
\textit{CU-Concerto-2026} с лоссом
\begin{equation}
\mathcal{L}_{\text{CNN}} = -\frac{1}{T}\sum_{t=1}^{T}
\sum_{k=1}^{88} y_{t,k}\,\log P_{\text{CNN}}(k\mid o_t),
\eqlab{cnn-loss}
\end{equation}
где $y_{t,k}\in\{0,1\}$ обозначает наличие $k$-й ноты в момент~$t$.
Оптимизатор Adam с шагом $3\!\times\!10^{-4}$, batch-size~$64$,
обучение продолжалось до сходимости валидационной F1-метрики
($24$ эпохи). Полученное распределение $P_{\text{CNN}}(\cdot\mid o_t)$
подставляется в~\Eqref{forward-recursion} вместо гауссовой эмиссии:
$b_i(o_t)\equiv P_{\text{CNN}}(p_i\mid o_t)$. Это снижает ошибку
выравнивания на полифонических фрагментах в~$1{,}6$ раза по
сравнению с базовой гауссовой моделью.

Граничные случаи отдельно классифицированы и обрабатываются явными
правилами. Долгие ферматы и педальные удержания проявляются как
большие значения $\alpha_t(\hat\varphi)$ при отсутствии новых
сильных onset-кандидатов; в этом случае счётчик длительности
$d$ растёт без перехода на следующее состояние, а темп оркестра
плавно замедляется до $0{,}7$ номинального. Резкие ритардандо в
каденции выявляются по знаковому изменению производной темпа в
окне $250$\,мс и инициируют включение специального профиля растяжения
с увеличенной чувствительностью. Полные пропуски тактов
обрабатываются процедурой повторной синхронизации, описанной
в~\Secref{realtime}: HMM переинициализируется на участке
$\pm 4$ ноты от наилучшей текущей оценки.

Орнаменты (форшлаги, мелизмы) представляют отдельную сложность:
их быстрые ноты не предусмотрены партитурой и могут вызывать
ложные продвижения HMM. В нашей реализации введён орнаментный
фильтр, основанный на оценке плотности onset’ов: если в окне
$60$\,мс зафиксировано более трёх onset’ов с убывающей энергией,
последовательность считается орнаментом и не передаётся трекеру.
Аналогичная эвристика используется для подавления случайных
повторений нот при двойных нажатиях.

Полный перечень рассмотренных граничных случаев и применяемых правил
сведён в \Tabref{corner-cases}.

\par\addvspace{2pt}
{\centering
\refstepcounter{table}\tablab{corner-cases}%
\footnotesize
\begin{tabular}{@{}p{0.42\linewidth}p{0.5\linewidth}@{}}
  \toprule
  Граничный случай                       & Правило обработки \\
  \midrule
  Фермата / удержание педали             & замораживание перехода, темп до $0{,}7$\,$\nu_0$ \\
  Резкое ритардандо в каденции           & профиль повышенной чувствительности RL \\
  Пропуск $\le 2$ нот                    & разрешённый skip-переход в~\Eqref{transition-matrix} \\
  Пропуск целого такта                   & повторная синхронизация HMM в окне $\pm 4$ ноты \\
  Случайный повтор ноты                  & фильтр двойных нажатий ($<25$\,мс) \\
  Орнамент / мелизм                      & окно $60$\,мс по плотности onset’ов \\
  Фальшивая нота полутоном               & принимается, $b_i(o_t)$ рассчитывается CNN \\
  Октавная ошибка ввода                  & замена кандидата на ближайшую к $p_{\hat\varphi}$ октаву \\
  \bottomrule
\end{tabular}\par
\vspace{2pt}
\tablename~\thetable. Перечень граничных случаев и эвристик их
обработки. Здесь $\nu_0$~--- номинальный темп фрагмента партитуры.\par}
\addvspace{4pt}

\sect{Опережающее предсказание темпа на основе RL}{rl}
%
Базовый трекер выдаёт оценку текущей позиции $\hat\varphi(t)$, однако
оркестровый рендерер, как показано в~\Secref{realtime}, нуждается в
\textit{будущем} темпе на горизонте $200$--$500$\,мс, чтобы успеть
синтезировать следующий блок без задержки. Простое экспоненциальное
сглаживание прошлых оценок плохо справляется с rubato: оно
переоценивает ускорения и недооценивает замедления, особенно в
кадансовых построениях. Мы переформулируем задачу предсказания темпа
как последовательную задачу принятия решений и обучаем для неё
отдельный модуль на основе обучения с подкреплением. Это
центральный методологический вклад настоящей работы и единственная
часть конвейера, в которой используется обучаемая на этапе работы
система.

Постановка опирается на дискретный марковский процесс принятия
решений $\mathcal{M}=(\mathcal{S},\mathcal{A},\mathcal{T},r,\gamma)$,
тик которого совпадает с тиком HMM (один шаг каждые $10$\,мс).
Состояние агента
\begin{equation}
s_t = \bigl(\alpha_t,\,\tau_{t-K:t},\,e_{t-K:t},\,
              \hat\varphi(t)\bigr)
\eqlab{rl-state}
\end{equation}
объединяет апостериорное распределение HMM, $K=20$ последних оценок
темпа, $K$ значений ошибки эмиссии (как косвенный индикатор
эмоционального напряжения исполнителя) и текущую позицию в партитуре.
Действие $a_t\in[0{,}5;\,2{,}0]$~--- мультипликативный темповый
коэффициент, применяемый рендерером в окне следующих $500$\,мс:
\begin{equation}
\nu_{t+1:t+500\,\text{мс}} = a_t\cdot\nu_0,
\eqlab{rl-action}
\end{equation}
где $\nu_0$~--- номинальный темп партитуры. Вознаграждение
\begin{equation}
\eqlab{rl-reward}
\begin{aligned}
r_t = \,&-\,\bigl|\,t^{\text{render}}_t - t^{\text{perf}}_t\,\bigr|
       \,-\, \lambda\,\mathcal{L}_{\text{align}}(t) \\
       &-\, \mu\,(a_t - a_{t-1})^2
\end{aligned}
\end{equation}
сочетает три компонента: рассогласование между моментом выпуска
оркестрового сэмпла и фактическим onset солиста, унаследованную потерю
выравнивания базового трекера и штраф за резкие изменения темпа,
обеспечивающий гладкость управления. Гиперпараметры $\lambda$ и $\mu$
подбираются на валидационной выборке, дисконт $\gamma=0{,}95$
соответствует горизонту $500$\,мс при шаге $10$\,мс.

\par\addvspace{2pt}
{\centering
\refstepcounter{figure}\figlab{rl-agent}%
\resizebox{\columnwidth}{!}{% =====================================================================
% TikZ figure: RL anticipation agent (BLACK & WHITE)
% Compact layout: 3 boxes top, 2 boxes bottom; fits into one column.
% Used in: sections/07-rl-anticipation.tex
% =====================================================================
\begin{tikzpicture}[
    node distance = 4mm and 3mm,
    every node/.append style={font=\scriptsize},
    block/.style ={rectangle, draw=black, thick, rounded corners=2pt,
                   minimum height=9mm, minimum width=20mm,
                   align=center, font=\scriptsize},
    arrow/.style ={-Stealth, thick}]

  \node[block]                  (track)  {Базовый\\трекер};
  \node[block, right=of track]  (state)  {Состояние\\$s_t$};
  \node[block, right=of state]  (policy) {Политика\\$\pi_\theta$};

  \node[block, below=7mm of policy] (env)    {Среда\\(рендерер)};
  \node[block, left=of env]         (reward) {Награда\\$r_t$};

  \draw[arrow] (track)  -- (state);
  \draw[arrow] (state)  -- (policy);
  \draw[arrow] (policy) -- node[right, font=\tiny] {$a_t$} (env);
  \draw[arrow] (env)    -- (reward);
  \draw[arrow] (reward.west) -| (track.south);

\end{tikzpicture}
}\par
\vspace{2pt}
\footnotesize
\figurename~\thefigure. Схема взаимодействия RL-модуля с базовым
трекером и средой. Агент получает состояние $s_t$, выдаёт темповый
коэффициент $a_t$ на горизонте $200$--$500$\,мс и обновляется по
вознаграждению $r_t$.\par}
\addvspace{4pt}

Обучение проводится в два этапа. На первом этапе строится политика
имитации: для каждой записи из CU-Concerto-2026 (\Secref{datasets})
рассчитывается «учительская» траектория темпа $a^*_t$ как
сглаженная производная времени выровненной MIDI-последовательности.
Политика $\pi_\theta(s_t)$ обучается минимизацией
\begin{equation}
\mathcal{L}_{\text{BC}}(\theta) =
\Expect_{(s,a^*)\sim\mathcal{D}}
\bigl\|\pi_\theta(s) - a^*\bigr\|_2^2,
\eqlab{rl-bc-loss}
\end{equation}
что соответствует классическому behavioural cloning. Архитектура
$\pi_\theta$~--- двухслойный перцептрон со скрытыми размерностями
$128$ и $64$; для учёта временного контекста дополнительно
тестируется однослойный GRU с $64$ скрытыми единицами.

На втором этапе политика дообучается алгоритмом Proximal Policy
Optimization~\cite{Schulman2017PPO} c использованием вознаграждения
\eqref{eq:rl-reward}. Параметры обучения: коэффициент клиппинга
$\epsilon=0{,}2$, $\lambda_{\text{GAE}}=0{,}95$, размер мини-батча
$2048$ переходов, $32$ эпохи PPO на батч, оптимизатор Adam со
скоростью $3\!\times\!10^{-4}$. Эпизоды формируются как мини-фрагменты
по $16$ тактов, что ускоряет распределение награды и позволяет
эффективно использовать GPU. Среда обучения~--- симулятор солиста,
воспроизводящий записанные в CU-Concerto-2026 исполнения с случайными
возмущениями темпа в полосе $\pm 15$\,\%.

Сводка гиперпараметров приведена в \Tabref{rl-hparams}.

\par\addvspace{2pt}
{\centering
\refstepcounter{table}\tablab{rl-hparams}%
\footnotesize
\begin{tabular}{@{}lcc@{}}
  \toprule
  Гиперпараметр                           & Обозн.\,         & Значение \\
  \midrule
  Окно истории                            & $K$              & $20$ ($200$\,мс) \\
  Диапазон действия                       & $a$              & $[0{,}5;\,2{,}0]$ \\
  Дисконт                                 & $\gamma$         & $0{,}95$ \\
  Вес выравнивания                        & $\lambda$        & $0{,}3$ \\
  Вес гладкости                           & $\mu$            & $0{,}1$ \\
  Клиппинг PPO                            & $\epsilon$       & $0{,}2$ \\
  GAE $\lambda$                           & $\lambda_{\text{GAE}}$ & $0{,}95$ \\
  Скорость обучения                       & $\eta$           & $3\!\times\!10^{-4}$ \\
  \bottomrule
\end{tabular}\par
\vspace{2pt}
\tablename~\thetable. Гиперпараметры обучения RL-модуля
предсказания темпа.\par}
\addvspace{4pt}

В режиме инференса модуль запрашивается параллельно с обновлением
HMM на каждом тике, что обеспечивает накладной расход около $1$\,мс
для MLP-варианта политики. Выход политики ограничивается
$[0{,}5;\,2{,}0]$ во избежание численных артефактов, а в случае
выхода за этот диапазон в лог записывается событие деградации.
Подробное описание среды, словаря состояний и формы наград приводится
в~\Secref{appendix-rl}.

В качестве альтернативных направлений для развития подхода в работе
сохраняются две конструкции, предложенные на этапе планирования.
Первая~--- замена явного вознаграждения~\eqref{eq:rl-reward} на
парные предпочтения, собираемые от музыкантов консерватории, по
схеме~\cite{Christiano2017Preference}: это позволит обучать темповую
интерпретацию непосредственно по эстетическим оценкам
профессионалов. Вторая~--- формализация всего трекера как POMDP с
belief-состоянием, равным forward-распределению HMM, и обучение
агента, выбирающего действия из дискретного множества \{«продолжать
следить», «перейти на OLTW», «запросить ресинхронизацию»\}~\cite{Kaelbling1998POMDP}.
Обе конструкции выходят за рамки настоящей статьи и обозначены как
направления дальнейших исследований.

\sect{Эксперименты}{experiments}
%
Эксперименты направлены на проверку трёх гипотез: (i) гибридная схема
HMM+OLTW сохраняет точность выравнивания при наличии rubato и редких
ошибок исполнителя; (ii) свёрточный фронтенд снижает ошибку
выравнивания на полифоническом материале по сравнению с гауссовой
эмиссией; (iii) RL-модуль опережающего предсказания темпа уменьшает
эффективную задержку аккомпанемента без ухудшения точности.

Базовые конфигурации, с которыми проводилось сравнение, обозначены
далее как \textit{HMM} (только скрытая марковская модель с гауссовой
эмиссией и фиксированной полосной матрицей переходов), \textit{HMM-CNN}
(тот же тракт, но с обучаемой эмиссией из~\Secref{cnn-corner}) и
\textit{HMM+OLTW+RL} (полная предложенная схема). Все три конфигурации
работают в общем конвейере (\Secref{realtime}) с одинаковыми
параметрами очистки и одинаковым рендерером, что позволяет
сосредоточить сравнение на алгоритмических различиях.

Метрики качества подобраны в соответствии с принятой в литературе
практикой~\cite{Nakamura2015,Henkel2025}. Средняя абсолютная ошибка
выравнивания MAE измеряется в номерах нот и характеризует, насколько
далеко предсказанная позиция отличается от истинной по партитуре.
Метрика F1-score рассчитывается как гармоническое среднее точности и
полноты для onset-событий, выровненных с допуском $50$\,мс. Метрика
ReSync~--- доля исполнения, на которой потребовалась повторная
синхронизация. Эффективная задержка $L_{\text{eff}}$ измеряется
как медиана временного сдвига между отрисованным оркестровым сэмплом
и соответствующим onset солиста.

Сначала рассмотрим зависимость задержки аккомпанемента от горизонта
предсказания темпа. На \Figref{tempo-latency} видно, что без
опережающего модуля эффективная задержка монотонно возрастает с
ростом горизонта (поскольку рендерер вынужден буферизовать всё
больший интервал), тогда как для гибридной схемы HMM+OLTW+RL зависимость
немонотонна и достигает минимума в области $250$--$350$\,мс. Это
согласуется с проектным значением горизонта $\gamma=0{,}95$.

\par\addvspace{2pt}
{\centering
\refstepcounter{figure}\figlab{tempo-latency}%
\resizebox{\columnwidth}{!}{% =====================================================================
% TikZ figure: tempo-anticipation horizon vs effective accompaniment
% latency (BLACK & WHITE).  Used in sections/08-experiments.tex.
% Axis labels are placed near the centre of each axis (not at the
% arrow tips) so they do not collide with the arrowheads.
% =====================================================================
\begin{tikzpicture}[x=0.9cm,y=0.075cm]
  \draw[->,thick] (0,0) -- (8.2,0);
  \draw[->,thick] (0,0) -- (0,62);

  \foreach \x/\labl in {1/100,2/150,3/200,4/250,5/300,6/400,7/500}
    \draw (\x,0) -- (\x,-1.5) node[below,font=\scriptsize] {\labl};
  \foreach \y in {10,20,30,40,50,60}
    \draw (0,\y) -- (-0.12,\y) node[left,font=\scriptsize] {\y};

  \foreach \y in {10,20,30,40,50}
    \draw[dotted,black!50] (0,\y) -- (7.2,\y);

  \draw[thick]
    plot[smooth] coordinates
      {(1,44) (2,41) (3,37) (4,34) (5,31) (6,28) (7,26)};
  \draw[thick,dashed]
    plot[smooth] coordinates
      {(1,23) (2,21) (3,20) (4,19) (5,19) (6,20) (7,21)};

  \node[font=\scriptsize,anchor=west] at (4.4,46) {базовый трекер};
  \node[font=\scriptsize,anchor=west] at (4.4,16) {HMM+OLTW+RL};

  \node[font=\scriptsize,below=4mm] at (4,-3) {Горизонт, мс};
  \node[font=\scriptsize,rotate=90] at (-1.0,32) {Задержка, мс};
\end{tikzpicture}
}\par
\vspace{2pt}
\footnotesize
\figurename~\thefigure. Зависимость эффективной задержки
аккомпанемента от горизонта предсказания темпа для базового
трекера и предложенной гибридной схемы.\par}
\addvspace{4pt}

Далее измерена устойчивость трекера при увеличении вариативности
rubato (\Figref{error-vs-rubato}). Базовый HMM начинает терять точность
выравнивания при $\sigma_\tau \ge 6$\,\%, тогда как HMM+OLTW+RL
сохраняет приемлемую ошибку вплоть до $\sigma_\tau\approx 12$\,\%.
Это объясняется двумя факторами: возможностью переключения на OLTW в
зонах низкой уверенности HMM (\Secref{realtime}) и опережающим
управлением темпом (\Secref{rl}), которое сглаживает резкие
ускорения.

\par\addvspace{2pt}
{\centering
\refstepcounter{figure}\figlab{error-vs-rubato}%
\resizebox{\columnwidth}{!}{% =====================================================================
% TikZ figure: rubato strength vs alignment error (BLACK & WHITE).
% Used in sections/08-experiments.tex.
% Axis labels are placed near the centre of each axis (not at the
% arrow tips) so they do not collide with the arrowheads.
% =====================================================================
\begin{tikzpicture}[x=0.95cm,y=0.10cm]
  \draw[->,thick] (0,0) -- (8.0,0);
  \draw[->,thick] (0,0) -- (0,42);

  \foreach \x/\labl in {1/2,2/4,3/6,4/8,5/10,6/12,7/14}
    \draw (\x,0) -- (\x,-1.1) node[below,font=\scriptsize] {\labl};
  \foreach \y in {5,10,15,20,25,30,35,40}
    \draw (0,\y) -- (-0.11,\y) node[left,font=\scriptsize] {\y};

  \foreach \y in {10,20,30}
    \draw[dotted,black!50] (0,\y) -- (7.2,\y);

  \draw[thick]
    plot[smooth] coordinates
      {(1,7) (2,10) (3,14) (4,19) (5,25) (6,31) (7,36)};
  \draw[thick,dashed]
    plot[smooth] coordinates
      {(1,6) (2,8) (3,11) (4,15) (5,19) (6,24) (7,28)};

  \node[font=\scriptsize,anchor=west] at (4.4,33) {HMM};
  \node[font=\scriptsize,anchor=west] at (4.4,22) {HMM+OLTW+RL};

  \node[font=\scriptsize,below=4mm] at (4,-3) {Сила rubato, \%};
  \node[font=\scriptsize,rotate=90] at (-1.1,21) {Ошибка, мс};
\end{tikzpicture}
}\par
\vspace{2pt}
\footnotesize
\figurename~\thefigure. Зависимость ошибки выравнивания от силы
rubato. Базовая схема HMM теряет точность раньше, чем гибридная
схема HMM+OLTW+RL.\par}
\addvspace{4pt}

Сводные количественные результаты по всем трём конфигурациям
приведены в \Tabref{main-results}. Гибридная схема HMM+OLTW+RL
превосходит базовый HMM на тестовой части CU-Concerto-2026 по всем
четырём метрикам: MAE снижается с $0{,}19$ до $0{,}11$ ноты, F1
возрастает с $0{,}947$ до $0{,}979$, ReSync падает в $2{,}5$ раза, а
эффективная задержка $L_{\text{eff}}$ опускается с $34$ до $19$\,мс.
Замена гауссовой эмиссии на CNN (HMM-CNN) даёт промежуточный
результат, что подтверждает значимость каждого из двух нововведений.

\par\addvspace{2pt}
{\centering
\refstepcounter{table}\tablab{main-results}%
\footnotesize
\begin{tabular}{@{}lcccc@{}}
  \toprule
  Конфигурация       & MAE, нот & F1     & ReSync, \% & $L_{\text{eff}}$, мс \\
  \midrule
  HMM                & $0{,}19$ & $0{,}947$ & $5{,}3$  & $34$ \\
  HMM-CNN            & $0{,}14$ & $0{,}968$ & $3{,}1$  & $30$ \\
  HMM+OLTW+RL        & $0{,}11$ & $0{,}979$ & $2{,}1$  & $19$ \\
  \bottomrule
\end{tabular}\par
\vspace{2pt}
\tablename~\thetable. Сводные количественные результаты на
тестовой части датасета CU-Concerto-2026. Меньшие значения MAE,
ReSync и $L_{\text{eff}}$~--- лучше; большее значение F1~--- лучше.\par}
\addvspace{4pt}

Результаты по отдельным сценариям исполнения, включающим
естественное rubato, пропуски нот, орнаменты и педальные удержания,
приведены в \Tabref{scenarios}. Падение F1 у предложенной схемы
наблюдается только в самых тяжёлых режимах (комбинированный стресс-тест
и низкое отношение сигнал/шум аудио), но даже в них значение
остаётся выше $0{,}93$.

\par\addvspace{2pt}
{\centering
\refstepcounter{table}\tablab{scenarios}%
\footnotesize
\begin{tabular}{@{}p{0.45\linewidth}cccc@{}}
  \toprule
  Сценарий исполнения                          & $\sigma_\tau$, \% & MAE & F1 & ReSync, \% \\
  \midrule
  Ровный темп, без ошибок                      & 2  & $0{,}07$ & $0{,}992$ & $0{,}0$ \\
  Умеренное rubato, без пропусков              & 6  & $0{,}09$ & $0{,}988$ & $0{,}0$ \\
  Сильное rubato в каденции                    & 11 & $0{,}12$ & $0{,}981$ & $0{,}7$ \\
  Случайные ускорения на переходах             & 8  & $0{,}11$ & $0{,}984$ & $0{,}5$ \\
  Резкое замедление перед кульминацией         & 10 & $0{,}13$ & $0{,}979$ & $0{,}9$ \\
  Единичный пропуск ноты                       & 4  & $0{,}10$ & $0{,}986$ & $1{,}1$ \\
  Серия из двух пропусков                      & 5  & $0{,}14$ & $0{,}975$ & $2{,}4$ \\
  Ложное повторение короткого мотива           & 7  & $0{,}16$ & $0{,}969$ & $3{,}1$ \\
  Орнамент с лишними быстрыми нотами           & 6  & $0{,}15$ & $0{,}972$ & $2{,}7$ \\
  Переход с педалью и размытым onset           & 7  & $0{,}18$ & $0{,}964$ & $4{,}8$ \\
  Низкое отношение сигнал/шум (аудио)          & 8  & $0{,}21$ & $0{,}951$ & $6{,}5$ \\
  Комбинированный стресс-тест                  & 12 & $0{,}27$ & $0{,}932$ & $9{,}8$ \\
  \bottomrule
\end{tabular}\par
\vspace{2pt}
\tablename~\thetable. Метрики по отдельным сценариям исполнения
для предложенной схемы HMM+OLTW+RL на CU-Concerto-2026.\par}
\addvspace{4pt}

Дополнительно проведена субъективная оценка с участием шести
студентов-консерваторов, не участвовавших в записи обучающего
датасета. Каждому испытуемому предлагалось сыграть одну из частей
концерта Чайковского с виртуальным аккомпанементом по двум
конфигурациям (HMM и HMM+OLTW+RL) в случайном порядке без указания,
какая из них активна. По завершении испытуемый оценивал три аспекта
по пятибалльной шкале: естественность темпового следования,
ощущение «живой реакции» оркестра и общее удобство репетиции.
Средние оценки HMM+OLTW+RL по всем трём аспектам выше базовой
конфигурации на $1{,}1$ балла, что согласуется с количественными
результатами и подтверждает практическую значимость
предложенной схемы.

\sect{Заключение}{conclusion}
%
В настоящей работе представлено единое математическое описание
моделей Hidden Markov Model и On-line Time Warping для задачи
отслеживания партитуры в реальном времени, а также архитектура
гибридного конвейера, объединяющего обе модели и снабжённого модулем
обучения с подкреплением для опережающего предсказания темпа на
горизонте $200$--$500$\,мс. Конвейер включает в себя обучаемый
свёрточный фронтенд, заменяющий гауссову функцию эмиссии HMM при
работе с аудиовходом, и систематический набор правил для обработки
граничных случаев~--- ферматы, орнаменты, пропуски, повторы и
октавные ошибки. Эксперименты на синтетическом датасете в стиле
прелюдий Шопена, на открытых корпусах MAPS и MSMD, а также на
собственном датасете \textit{CU-Concerto-2026}, собранном в
Центральном Университете, показали, что предложенная схема
снижает эффективную задержку аккомпанемента с $34$ до $19$\,мс и
улучшает F1-метрику с $0{,}947$ до $0{,}979$ по сравнению с базовой
HMM-конфигурацией.

Главный методологический вклад работы~--- демонстрация того, что
обучение с подкреплением полезно не как замена классического трекера,
а как отдельный опережающий слой, действующий поверх него. Такой
подход сохраняет вероятностную интерпретируемость скрытых марковских
моделей, преимущества OLTW по устойчивости к рассогласованию и
одновременно даёт инструмент адаптации темпа под индивидуальный
стиль исполнителя.

Дальнейшие исследования планируется развивать в трёх направлениях.
Первое~--- замена явного вознаграждения~\eqref{eq:rl-reward} на парные
предпочтения, собираемые от музыкантов консерватории, по
схеме~\cite{Christiano2017Preference}: это позволит включить в
обучение эстетические критерии, не сводимые к одной численной метрике.
Второе~--- формализация всего трекера как POMDP над
belief-состоянием HMM с агентом, выбирающим действия из дискретного
множества альтернатив~\cite{Kaelbling1998POMDP}. Третье~--- расширение
датасета \textit{CU-Concerto-2026} и репертуара поддерживаемых
произведений, а также пилотное внедрение в учебный процесс
Центрального Университета и партнёрских консерваторий.


\sect{Благодарности}{acknowledgments}
%
Авторы благодарят менторов мастерской MusicTech Центрального
Университета и Т-Банка за постановку задачи и техническую
поддержку, а также студентов-консерваторов, согласившихся принять
участие в записи тестового датасета и в субъективной оценке
аккомпанемента.

\bibliographystyle{apsrev4-2}
\bibliography{references}

% =====================================================================
% Appendices start on a NEW page (as in пример/main.pdf), with a
% one-column "ПРИЛОЖЕНИЯ" banner, then individual sub-appendices in
% the usual two-column layout.
% =====================================================================
\clearpage
\onecolumngrid
\vspace{0.7\baselineskip}
\begingroup
  \centering
  \normalsize\bfseries
  ПРИЛОЖЕНИЯ\par
\endgroup
\vspace{0.25\baselineskip}
\twocolumngrid

\sect{Приложение А. Алгоритмы Forward и Viterbi}{appendix-hmm}
%
В настоящем приложении приведены полные рекуррентные формулы и
псевдокод классических алгоритмов Forward и Viterbi для скрытой
марковской модели, используемой в~\Secref{background}. Для
самостоятельности изложения здесь повторяется обозначение
$\lambda=(A,B,\pi)$, где $A=(A_{ij})$~--- матрица переходов согласно
\eqref{eq:transition-matrix}, $b_i(o)$~--- эмиссионная плотность,
$\pi=(\pi_i)$~--- начальное распределение.

Алгоритм Forward вычисляет полную вероятность наблюдаемой
последовательности $O=(o_1,\dots,o_T)$ при модели $\lambda$ через
вспомогательные величины
$\alpha_t(i) = \Prb(o_1,\dots,o_t,\,s_t=i\mid\lambda)$. Рекуррентная
формула имеет вид
\begin{equation}
\alpha_t(j) = b_j(o_t)\sum_{i=1}^{N}\alpha_{t-1}(i)\,A_{ij},
\eqlab{forward-app}
\end{equation}
с инициализацией $\alpha_1(j)=\pi_j\,b_j(o_1)$. Полная вероятность
наблюдений равна $\Prb(O\mid\lambda)=\sum_{j=1}^{N}\alpha_T(j)$. На
практике для длинных последовательностей применяется покадровая
нормировка $\tilde\alpha_t(j)=\alpha_t(j)\bigl/\sum_k\alpha_t(k)$,
что устраняет численное переполнение и одновременно даёт
апостериорную оценку $\Prb(s_t=j\mid o_1,\dots,o_t)$, используемую
для онлайн-выдачи $\hat\varphi(t)=\argmaxop_j\tilde\alpha_t(j)$.

\par\addvspace{2pt}
{\centering
\begin{minipage}{0.97\linewidth}
\footnotesize
\textbf{Алгоритм 1.} Forward-рекурсия для онлайн-трекера. \\[1pt]
\rule{\linewidth}{0.4pt}\vspace{1pt}\\
\textbf{Вход:} партитура $S=(p_1,\dots,p_N)$, матрица $A$,
эмиссионный параметр $\sigma$. \\
1: $\alpha\gets(1,0,\dots,0)\in\Real^N$ \\
2: \textbf{для} каждой поступающей ноты $o$ \textbf{выполнить} \\
3: \quad \textbf{для} $j\gets 1,\dots,N$ \textbf{выполнить} \\
4: \quad\quad $b_j\gets\mathcal{N}(o;\,p_j,\,\sigma^2)$ \\
5: \quad\quad $\alpha'_j\gets b_j\sum_{i=1}^{N}\alpha_i\,A_{ij}$ \\
6: \quad \textbf{конец для} \\
7: \quad $\alpha\gets\alpha'\bigl/\sum_k\alpha'_k$ \\
8: \quad \textbf{выдать} $\hat\varphi\gets\argmaxop_j\alpha_j$ \\
9: \textbf{конец для} \\
\rule{\linewidth}{0.4pt}
\end{minipage}\par}
\addvspace{4pt}

Алгоритм Viterbi возвращает не апостериорное распределение, а одну
наиболее вероятную траекторию состояний $s_1^*,\dots,s_T^*$. Он
оперирует величинами
$\delta_t(i) = \max_{s_1,\dots,s_{t-1}}
\Prb(s_1,\dots,s_t=i,\,o_1,\dots,o_t\mid\lambda)$
и обратными указателями $\psi_t(i)$:
\begin{equation}
\eqlab{viterbi-delta}
\delta_t(j) = b_j(o_t)\,\max_{i}\bigl[\delta_{t-1}(i)\,A_{ij}\bigr],
\end{equation}
\begin{equation}
\eqlab{viterbi-psi}
\psi_t(j) = \argmaxop_{i}\bigl[\delta_{t-1}(i)\,A_{ij}\bigr],
\end{equation}
с инициализацией $\delta_1(j)=\pi_j\,b_j(o_1)$ и $\psi_1(j)=0$.
Оптимальная траектория восстанавливается обратным проходом:
$s_T^*=\argmaxop_j\delta_T(j)$, $s_{t-1}^*=\psi_t(s_t^*)$.

В нашем трекере Viterbi используется в \textit{режиме отложенного
анализа}: после исполнения каждого фрагмента партитуры (например,
после двойной чёрточки) выполняется обратный проход по
накопленным $\delta$ и $\psi$, что позволяет уточнить выровнянную
последовательность для дальнейшего обучения CNN-эмиссии и
обновления статистик переходов на следующих исполнениях.

Сложность одного шага обоих алгоритмов составляет $O(N)$ за счёт
полосной структуры матрицы $A$ из~\eqref{eq:transition-matrix}:
вместо суммирования по всем состояниям достаточно обработать три
ненулевых перехода. Память~--- $O(N)$ для Forward и $O(N T)$ для
Viterbi (требуется хранить указатели $\psi_t$ за весь интервал
анализа).

\sect{Приложение Б. Псевдокод и анализ On-line Time Warping}{appendix-oltw}
%
В настоящем приложении приведены подробности On-line Time Warping,
используемого как резервный трекер в~\Secref{realtime}. Алгоритм
основан на каузальной версии классической DTW-рекурсии
\eqref{eq:dtw-recursion}, в которой матрица накопленных стоимостей
$D$ заполняется не целиком, а столбец за столбцом по мере поступления
новых наблюдений.

Состояние OLTW описывается тремя величинами: текущей строкой
$i_{\text{cur}}$ (последняя обработанная позиция партитуры), текущим
столбцом $j_{\text{cur}}$ (число принятых наблюдений) и счётчиком
\textit{RunCount}, ограничивающим число последовательных шагов в одну
сторону. Параметр $R_{\max}$ устанавливает верхнюю границу для
RunCount; в нашей реализации $R_{\max}=3$.

\par\addvspace{2pt}
{\centering
\begin{minipage}{0.97\linewidth}
\footnotesize
\textbf{Алгоритм 2.} OLTW в каузальном режиме. \\[1pt]
\rule{\linewidth}{0.4pt}\vspace{1pt}\\
\textbf{Вход:} партитура $S=(p_1,\dots,p_N)$, ширина окна $W$,
ограничение $R_{\max}$. \\
1: $i_{\text{cur}}\gets 1,\,j_{\text{cur}}\gets 1$,
   \,RunCount $\gets 0$, \,direction $\gets 0$ \\
2: $D[1,1]\gets d(s_1,q_1)$ \\
3: \textbf{для} каждой поступающей ноты $q_j$ \textbf{выполнить} \\
4: \quad обновить столбец $j$ матрицы $D$ в окне
   $[\,i_{\text{cur}}-W,\,i_{\text{cur}}+W\,]$
   по~\eqref{eq:dtw-recursion} \\
5: \quad $(i^*,j^*)\gets\argminop\bigl\{D[i,j]:
   \;|\,i-i_{\text{cur}}|\le 1,\,|\,j-j_{\text{cur}}|\le 1\bigr\}$ \\
6: \quad \textbf{если} $(i^*,j^*)$ продолжает текущее
   direction \textbf{то} \\
7: \quad\quad RunCount $\gets$ RunCount $+\,1$ \\
8: \quad \textbf{иначе} \\
9: \quad\quad RunCount $\gets 1$;\,\, direction $\gets$
   шаг$(i_{\text{cur}}\!\to\!i^*,\,j_{\text{cur}}\!\to\!j^*)$ \\
10: \quad \textbf{конец если} \\
11: \quad \textbf{если} RunCount $\ge R_{\max}$ \textbf{то} \\
12: \quad\quad принудительно сменить direction
   на ортогональное, RunCount $\gets 0$ \\
13: \quad \textbf{конец если} \\
14: \quad $i_{\text{cur}}\gets i^*,\;j_{\text{cur}}\gets j^*$ \\
15: \quad \textbf{выдать} $\hat\varphi(j)\gets i_{\text{cur}}$ \\
16: \textbf{конец для} \\
\rule{\linewidth}{0.4pt}
\end{minipage}\par}
\addvspace{4pt}

Ширина окна $W$ задаёт компромисс между робастностью и
вычислительной стоимостью: при $W=N$ алгоритм совпадает с полной
DTW (но теряет онлайн-свойство), при $W=1$ вырождается в жадное
сравнение соседних позиций. В наших экспериментах
$W=\lceil 0{,}1\,N\rceil$, что соответствует допуску по позиции в
полосе примерно $10$\,\% длины партитуры. Сложность одного шага
составляет $O(W)$, суммарная сложность для исполнения длиной
$M$~--- $O(M W)$, что при $W\ll N$ существенно меньше
$O(NM)$ полной DTW.

Геометрический смысл механизма \textit{RunCount} проиллюстрирован
ниже. Без него каузальная DTW могла бы непрерывно двигаться по
одной строке (горизонтально), если бы это локально минимизировало
накопленную стоимость, и оказалась бы заблокированной в
произвольной позиции партитуры. Принудительное переключение
направления после $R_{\max}$ шагов гарантирует продвижение по
обеим осям и исключает «застревание». Эмпирически $R_{\max}=3$
оказывается оптимальным: меньшие значения дают рваное
выравнивание на тонкой структуре звукоряда, бо́льшие~--- допускают
застревание на ферматах.

В сочетании с HMM (см.~\Secref{realtime}) OLTW служит резервным
трекером, активируемым при низкой уверенности forward-распределения.
В этом режиме OLTW инициализируется текущей наилучшей оценкой
$\hat\varphi$ HMM и заменяет её на интервале до восстановления
уверенности. После восстановления HMM продолжается с
$\alpha$-вектором, переинициализированным в окне $\pm 4$ ноты
вокруг последнего значения OLTW.

\sect{Приложение В. Среда RL и подробности реализации}{appendix-rl}
%
В настоящем приложении детально описана среда, на которой обучается и
оценивается RL-модуль предсказания темпа из~\Secref{rl}. Среда
реализована поверх симулятора, читающего размеченные исполнения из
датасета \textit{CU-Concerto-2026}, и снабжена API, совместимым с
библиотекой \texttt{gymnasium}. Один эпизод соответствует одному
шестнадцатитактному фрагменту партитуры, что даёт длину эпизода
порядка $200$ шагов агента при базовом темпе $\nu_0=120$\,bpm.

Пространство состояний $\mathcal{S}\subset\Real^{2N+2K+1}$. Координаты
включают: forward-распределение $\alpha_t\in\Real^N$, текущую
позицию $\hat\varphi(t)/N$, $K=20$ последних оценок темпа
$\tau_{t-K:t}$ (нормированных к $\nu_0$), $K$ последних значений
эмиссионной ошибки $e_{t-K:t}=-\log b_{\hat\varphi}(o_t)$ и
относительное время с начала эпизода. Для уменьшения размерности
$\alpha_t$ применяется проекция в спектральные коэффициенты
дискретного косинусного преобразования первого рода
$\alpha_t\mapsto\hat\alpha_t\in\Real^{16}$. Эмпирически такая
проекция сохраняет более $97$\,\% энергии распределения и снижает
требования к политике почти на порядок.

Пространство действий~--- одномерный отрезок $\mathcal{A}=[0{,}5;\,2{,}0]$.
Политика моделируется усечённым гауссовским распределением
$\pi_\theta(a\mid s)=\mathcal{N}_{\text{trunc}}(\mu_\theta(s),
\sigma_\theta^2)\bigr|_{[0{,}5;\,2{,}0]}$, что обеспечивает
гладкое исследование без резких выходов из физически разумного
диапазона.

Функция вознаграждения определена~\eqref{eq:rl-reward} и состоит из
трёх компонентов. Первый~--- штраф за рассогласование
$|t^{\text{render}}_t-t^{\text{perf}}_t|$~--- измеряется в
миллисекундах, нормирован на $50$\,мс (поскольку именно это значение
является порогом восприятия рассинхронизации в музыкальном
ансамбле, см.~\Secref{introduction}). Второй компонент~---
$\mathcal{L}_{\text{align}}(t)$~--- равен расстоянию по партитуре
между ground-truth и текущей оценкой $\hat\varphi$, нормированному
на длину фрагмента. Третий компонент~--- штраф гладкости
$(a_t-a_{t-1})^2$~--- ограничивает резкие переключения темпа,
которые акустически воспринимаются как «дёрганье» оркестра.

Симулятор работает в двух режимах. В \textit{train-режиме} к
номинальному темпу записанного исполнения применяются случайные
возмущения, эмулирующие индивидуальное rubato: умножение
длительностей на гауссовский шум $\mathcal{N}(1,\sigma_\tau^2)$,
случайные локальные ритардандо и ачелерандо в окне восьми тактов и
вставка пропусков с вероятностью $0{,}5$\,\% на ноту. В
\textit{eval-режиме} симулятор воспроизводит ровно записанное
исполнение без возмущений; этот режим используется для
итоговой оценки на отложенной выборке.

\par\addvspace{2pt}
{\centering
\refstepcounter{table}\tablab{rl-env}%
\footnotesize
\begin{tabular}{@{}lc@{}}
  \toprule
  Параметр среды                           & Значение \\
  \midrule
  Длина эпизода                            & $16$ тактов ($\sim 200$ шагов) \\
  Шаг агента                               & $10$\,мс \\
  Размерность состояния                    & $73$ ($16+20+20+1+16$) \\
  Размерность действия                     & $1$ \\
  Возмущение темпа в train-режиме          & $\sigma_\tau\in[0{,}02;\,0{,}12]$ \\
  Вероятность пропуска ноты                & $5\!\times\!10^{-3}$ \\
  Вероятность повтора ноты                 & $2\!\times\!10^{-3}$ \\
  Норма для штрафа рассогласования         & $50$\,мс \\
  Норма для штрафа выравнивания            & длина фрагмента \\
  \bottomrule
\end{tabular}\par
\vspace{2pt}
\tablename~\thetable. Параметры симулятора, реализующего среду
для обучения и оценки RL-модуля предсказания темпа.\par}
\addvspace{4pt}

Реализация написана на Python~3.11 с использованием библиотеки
\texttt{stable-baselines3} в части PPO и \texttt{torch} версии $2{,}3$
в части политики. Все вычисления, не связанные с обучением, проводятся
на CPU; обучение PPO~--- на одной GPU NVIDIA~RTX~3060. Полная сходимость
достигается за $\sim 12$ часов при $5\!\times\!10^{6}$ шагах среды;
behavioural cloning требует около $25$ минут на тех же исходных
данных. Соответствующие веса политики сохраняются в формате
\texttt{state\_dict} и загружаются в продакшн-конвейер
(\Secref{realtime}) как обычный PyTorch-чекпойнт.


\end{document}
